\DocumentMetadata{
  lang        = en,
	pdfversion  = 1.7,
  pdfstandard = ua-1,
%  pdfstandard = a-4f, %or a-4
  tagging=on,
	tagging-setup={math/mathml/luamml/load=false}
%  tagging-setup={math/setup=mathml-SE}
}
%\tagpdfsetup{role/map-tags=pdf}
%\tagpdfsetup{math/mathml/luamml/load=false}
\tagpdfsetup{table/header-rows={1}}

\documentclass[10pt]{article}

\usepackage[margin=1in]{geometry}

\author{Math 103: College Algebra}
\date{}
%\author{Dr. Masaros}
\pagestyle{empty}
\pagenumbering{gobble}

\usepackage{xparse}
%\usepackage{enumerate}
%\usepackage{mdwlist}
%\usepackage[inline]{enumitem}
\usepackage{enumitem}
\usepackage{tabto}

%\usepackage{amsmath,amssymb,amsthm}
%\usepackage{amsmath}
%\usepackage{unicode-math}
\usepackage{amsmath,amssymb,amsthm}
%\usepackage{amsmath}

\usepackage{textcomp}
\newcommand{\cent}{{\textcent}}

%Tagging is incompatible with titlesec
%\usepackage[compact,tiny]{titlesec}
%\usepackage{minipage}
\makeatletter
\renewcommand\section{%
  \@startsection{section}% name
    {1}% level
    {0pt}% indent
    {-1ex plus -.3ex minus -.06ex}% beforeskip
    {1ex plus .1ex}% afterskip
    {\bfseries}% style
  }%
\makeatother

\usepackage{titling}
\setlength{\droptitle}{-.5in}
\pretitle{\centering\noindent\LARGE}
\posttitle{\par}
\preauthor{\centering\noindent\large}
\postauthor{\par}
\predate{\centering\noindent\large}
\postdate{\par}

\newcommand{\NameLine}{\noindent Full Name: \rule{\linewidth-\widthof{Full Name: }}{.4pt}}

\newcommand{\MakeLine}[1][1in]{\rule{#1}{.4pt}}
\newcommand{\MakeLineFull}[1]{\noindent #1: \rule{\linewidth-\widthof{#1: }}{.4pt}}
\newcommand{\MakeLineSpace}[1][1in]{\rule{0pt}{0.34375in}\rule{#1}{.4pt}}
\newcommand{\MakeLineFullSpace}[1]{\noindent #1: \rule{0pt}{0.34375in}\rule{\linewidth-\widthof{#1: }}{.4pt}}

\usepackage{calc}
\newlength{\WorkSpace}
\setlength{\WorkSpace}{.25\textheight}
\newcommand{\MakeSpace}[1][\WorkSpace]{\rule{0pt}{#1}}
\newcommand{\itemS}[1][\WorkSpace]{\item\MakeSpace[#1]}
\newcommand{\itemCB}{\vfill{ }\columnbreak\item}
\newcommand{\itemC}[1][]{\clearpage\item}

\newcommand{\Hspace}[1]{\rule{#1}{0pt}}

\usepackage{tabularx}
\newcolumntype{C}{>{\centering\arraybackslash}X}
\newcolumntype{R}{>{\hfill\arraybackslash}X}

\newcommand{\ph}{\phantom{1}}
\newcommand{\phc}{\phantom{1,}}
\newcommand{\phm}{\phantom{-}}

\usepackage{xparse}

\usepackage{tikz}
\usetikzlibrary{decorations.markings}
\usepgflibrary{arrows}
\usetikzlibrary{intersections}
\usetikzlibrary{calc}

\tikzset{>=angle 90}
\tikzset{x=.5in,y=.5in}
\tikzset{every picture/.append style={artifact}}

\newcommand{\MakeBlank}[1][10em]{\rule{#1}{.4pt}}
\newcommand{\FullLine}{\\[3ex]\MakeBlank[\linewidth-1pt]}

\newenvironment{cpic}[1][]{%
\begin{center}
\begin{tikzpicture}[#1]
}{%
\end{tikzpicture}
\end{center}
}

\newenvironment{ilpic}[1][x=0.2in,y=0.2in]{%
\begin{tikzpicture}[#1,baseline={([yshift=-2ex]current bounding box.north)}]
}{%
\end{tikzpicture}
}

%\usepackage[scr=boondoxo]{mathalfa}
%\usepackage[margin=1in]{geometry}


\usepackage{multienum}

\newcommand{\Line}[1]{\ensuremath{\overleftrightarrow{\;#1\;}}}
\newcommand{\Ray}[1]{\ensuremath{\overrightarrow{#1\;}}}
\newcommand{\Segment}[1]{\ensuremath{\overline{#1}}}

\renewcommand{\ell}{\ensuremath{\mathscr{l}}}
\newcommand{\emm}{\ensuremath{\mathscr{m}}}
\newcommand{\enn}{\ensuremath{\mathscr{n}}}
\newcommand{\pee}{\ensuremath{\mathscr{p}}}

\newcommand{\Angle}{\angle}

\newcommand{\Mathscr}[1]{\ensuremath{\mathscr{#1}}}

\usepackage{multicol}

\usepackage{setspace}

\newcounter{enumSave}
\newcommand{\saveenum}{\setcounter{enumSave}{\value{enumi}}}
\newcommand{\loadenum}{\setcounter{enumi}{\value{enumSave}}}

\newcommand{\savemenum}{\setcounter{enumSave}{\value{multienumi}}}
\newcommand{\loadmenum}{\setcounter{multienumi}{\value{enumSave}}}

\newcommand{\Norm}[1]{\left|#1\right|}

\newenvironment{senum}{\begin{enumerate}\loadenum}{\saveenum\end{enumerate}}
\newenvironment{smenum}{\begin{multienumerate}\loadmenum}{\savemenum\end{multienumerate}}

\newenvironment{nmenum}{\setlength\hsize{\linewidth}\begin{multienumerate}\renewcommand{\labelenumi}{\addtocounter{multienumi}{1}\alph{multienumi}.}}{\end{multienumerate}}
\newenvironment{Nmenum}{\setlength\hsize{\linewidth}\begin{multienumerate}\renewcommand{\labelenumi}{\addtocounter{multienumi}{1}\Alph{multienumi}.}}{\end{multienumerate}}

\newcommand{\pentagon}{%
\begin{tikzpicture}[x=1ex,y=1ex]%
\draw (90:1) -- (90+72:1) -- (90+72+72:1) -- (90-72-72:1) -- (90-72:1) -- cycle;%
\end{tikzpicture}%
}

\usepackage{environ}

\usepackage{forloop}
\newcounter{ThisVersion}
\newcounter{SavePage}
\newcounter{SaveSection}

\NewEnviron{MultiVersion}[1]{
\forLoop{1}{#1}{ThisVersion}{%
\setcounter{SavePage}{\value{page}}%
\setcounter{SaveSection}{\value{section}}
\BODY%
\clearpage
\setcounter{page}{\value{SavePage}}%
\setcounter{section}{\value{SaveSection}}
}
}

\usepackage{ifthen}
\newcommand{\vitem}[2]{%
\ifthenelse{\value{ThisVersion}=#1}{\item #2}{}%
}

\newcommand{\QuadrilateralChoices}{%
\begin{nmenum}
\mitemxxxx
{Parallelogram}
{Rectangle}
{Rhombus}
{Square}
\end{nmenum}
}

\newcommand{\MultiLines}[1][3in]{%
\\\noindent\begin{tikzpicture}[y=.25in, ultra thin]
\draw[xstep=2\textwidth,ystep=1] (.25\textwidth,0) grid (1.25\textwidth,#1);
\end{tikzpicture}
}

\newcommand{\DrawProtractor}[1][blue!10!white]{
\filldraw[fill=#1] (20:4) ++(110:-.5) -- ++(110:.5) arc (20:200:4) -- ++(110:-.5) -- cycle;
\draw (20:3) -- (200:3);
\filldraw[fill=white] (30:3) arc (30:190:3) -- cycle;
\foreach \t in {0,10,...,180}
    {
		\draw (\t+20:4) -- (\t+20:3.75);
		\draw (\t+20:3.625) node[rotate=-70+\t]{$\t$};
		\draw (\t+20:3) -- (\t+20:3.25);
		\draw (200-\t:3.375) node[rotate=110-\t]{$\t$};
    }
\foreach \t in {0,10,...,170}
		\foreach \d in {1,2,...,9}
		    \draw (\t+\d+20:3.825) -- (\t+\d+20:4);
}

\newcounter{EqNum}
\setcounter{EqNum}{0}
\newcommand{\EqLbl}{\addtocounter{EqNum}{1}(\Alph{EqNum})}

\newenvironment{esys}{
\setcounter{EqNum}{0}
\displaystyle
\left\{\begin{array}{r@{\;}r@{\;}r@{\;}r@{\;}r@{\quad\EqLbl}}
}{
\end{array}\right.}

\newenvironment{fsys}{
\setcounter{EqNum}{0}
\displaystyle
\left\{\begin{array}{l@{\quad\EqLbl}}
}{
\end{array}\right.}

\NewDocumentCommand{\DrawGridSquare}{O{5} D(){1}}{
\draw[step=#2,black!50!white] (-#1-.9,-#1-.9) grid (#1+.9,#1+.9);
\draw[<->,very thick] (-#1-1,0) -- (#1+1,0) node[anchor=south]{$x$};
\draw[<->,very thick] (0,-#1-1) -- (0,#1+1) node[anchor=east]{$y$};
\foreach \x in {1,...,#1}
    {
    \draw[thick] (\x,.2) -- (\x,-.2) node[anchor=north,font=\scriptsize]{$\x$};
    \draw[thick] (-\x,.2) -- (-\x,-.2) node[anchor=north,font=\scriptsize]{$-\x\phm$};
    \draw[thick] (.2,\x) -- (-.2,\x) node[anchor=east,font=\scriptsize]{$\x$};
    \draw[thick] (.2,-\x) -- (-.2,-\x) node[anchor=east,font=\scriptsize]{$-\x$};
    }
}

\newcommand{\DrawGridSquareEO}[1][5]{
\draw[step=1,black!50!white] (-#1-.9,-#1-.9) grid (#1+.9,#1+.9);
\draw[<->,very thick] (-#1-1,0) -- (#1+1,0) node[anchor=south]{$x$};
\draw[<->,very thick] (0,-#1-1) -- (0,#1+1) node[anchor=east]{$y$};
\foreach \x in {2,4,...,#1}
    {
    \draw[thick] (\x,.2) -- (\x,-.2) node[anchor=north,font=\scriptsize]{$\x$};
    \draw[thick] (-\x,.2) -- (-\x,-.2) node[anchor=north,font=\scriptsize]{$-\x\phm$};
    \draw[thick] (.2,\x) -- (-.2,\x) node[anchor=east,font=\scriptsize]{$\x$};
    \draw[thick] (.2,-\x) -- (-.2,-\x) node[anchor=east,font=\scriptsize]{$-\x$};
    }
}

\newcounter{t1}
\newcounter{t2}

\newcommand{\DrawGrid}[5][1]{
\draw[step=1,black!50!white] (-#2-.9,-#4-.9) grid (#3+.9,#5+.9);
\draw[<->,very thick] (-#2-1,0) -- (#3+1,0) node[anchor=south]{$x$};
\draw[<->,very thick] (0,-#4-1) -- (0,#5+1) node[anchor=east]{$y$};
\setcounter{t1}{#1}
\setcounter{t2}{2*\value{t1}}
\ifthenelse{#3<\value{t1}}{}{%
\ifthenelse{#3<\value{t2}}{\draw[thick] (\arabic{t1},0.05in) -- (\arabic{t1},-0.05in) node[anchor=north,font=\scriptsize]{$\arabic{t1}$};}{
\foreach \x in {\arabic{t1},\arabic{t2},...,#3}
    \draw[thick] (\x,0.05in) -- (\x,-0.05in) node[anchor=north,font=\scriptsize]{$\x$};
}}
\ifthenelse{#2<\value{t1}}{}{%
\ifthenelse{#2<\value{t2}}{\draw[thick] (-\arabic{t1},0.05in) -- (-\arabic{t1},-0.05in) node[anchor=north,font=\scriptsize]{$-\arabic{t1}\phm$};}{
\foreach \x in {\arabic{t1},\arabic{t2},...,#2}
    \draw[thick] (-\x,0.05in) -- (-\x,-0.05in) node[anchor=north,font=\scriptsize]{$-\x\phm$};
}}
\ifthenelse{#5<\value{t1}}{}{%
\ifthenelse{#5<\value{t2}}{\draw[thick] (0.05in,\arabic{t1}) -- (-0.05in,\arabic{t1}) node[anchor=east,font=\scriptsize]{$\arabic{t1}$};}{
\foreach \x in {\arabic{t1},\arabic{t2},...,#5}
    \draw[thick] (0.05in,\x) -- (-0.05in,\x) node[anchor=east,font=\scriptsize]{$\x$};
}}
\ifthenelse{#5<\value{t1}}{}{%
\ifthenelse{#5<\value{t2}}{\draw[thick] (0.05in,-\arabic{t1}) -- (-0.05in,-\arabic{t1}) node[anchor=east,font=\scriptsize]{$-\arabic{t1}$};}{
\foreach \x in {\arabic{t1},\arabic{t2},...,#4}
    \draw[thick] (0.05in,-\x) -- (-0.05in,-\x) node[anchor=east,font=\scriptsize]{$-\x$};
}}
}

\NewDocumentCommand{\GridSquareFit}{O{5} D(){2.75in}}{%
\begin{ilpic}[x={#2/(2*#1+2)},y={#2/(2*#1+2)}]
\useasboundingbox (-#1-1,-#1-1) rectangle (#1+1,#1+1);
\DrawGridSquare[#1]
\end{ilpic}
}

\NewDocumentCommand{\GridSquareFitHalf}{O{5} D(){2.75in}}{%
\begin{ilpic}[x={#2/(2*#1+2)},y={#2/(2*#1+2)}]
\useasboundingbox (-#1-1,-#1-1) rectangle (#1+1,#1+1);
\DrawGridSquare[#1](0.5)
\end{ilpic}
}

\NewDocumentCommand{\GridSquareFitEO}{O{5} D(){2.75in}}{%
\begin{ilpic}[x={#2/(2*#1+2)},y={#2/(2*#1+2)}]
\useasboundingbox (-#1-1,-#1-1) rectangle (#1+1,#1+1);
\DrawGridSquareEO[#1]
\end{ilpic}
}

\NewDocumentEnvironment{GridSquareFitEOFilled}{O{5} D(){2.75in} +b}{%
\begin{ilpic}[x={#2/(2*#1+2)},y={#2/(2*#1+2)}]
\useasboundingbox (-#1-1,-#1-1) rectangle (#1+1,#1+1.5);
\DrawGridSquareEO[#1]
#3
\end{ilpic}
}{}

\NewDocumentEnvironment{GridFitFilledEO}{m m m m D(){2.75in} +b}{%
\begin{ilpic}[x={#5/(#2+#1+2)},y={#5/(#4+#3+2)}]
\useasboundingbox (-#1-1,-#3-1) rectangle (#2+1,#4+1.5);
\DrawGrid[2]{#1}{#2}{#3}{#4}
#6
\end{ilpic}
}{}

\NewDocumentEnvironment{GridFitFilled}{O{1} m m m m D(){2.75in} +b}{%
\begin{ilpic}[x={#6/(#3+#2+2)},y={#6/(#5+#4+2)}]
\useasboundingbox (-#2-1,-#4-1) rectangle (#3+1,#5+1.5);
\DrawGrid[#1]{#2}{#3}{#4}{#5}
#7
\end{ilpic}
}{}

\NewDocumentEnvironment{GridSquareFitFilled}{O{5} D(){2.75in} +b}{%
\begin{ilpic}[x={#2/(2*#1+2)},y={#2/(2*#1+2)}]
\useasboundingbox (-#1-1,-#1-1) rectangle (#1+1,#1+1.5);
\DrawGrid[1]{#1}{#1}{#1}{#1}
#3
\end{ilpic}%
}{}

\NewDocumentEnvironment{GridFilled}{O{0.2in} m m m m O{1} +b}{%
\begin{ilpic}[x=#1,y=#1]%
\useasboundingbox(-#2-1,-#4-1) rectangle (#3+1,#5+1);
\DrawGrid[#6]{#2}{#3}{#4}{#5}
#7
\end{ilpic}%
}

\NewDocumentEnvironment{GridSquareFilled}{O{0.2in} m O{1} +b}{%
\begin{GridFilled}[#1]{#2}{#2}{#2}{#2}[#3]%
#4
\end{GridFilled}%
}

\newcommand{\DrawGridSquareBlank}[1][5]{
\draw[step=1,ultra thin] (-#1-.9,-#1-.9) grid (#1+.9,#1+.9);
\draw[<->,very thick] (-#1-1,0) -- (#1+1,0) node[anchor=south]{$x$};
\draw[<->,very thick] (0,-#1-1) -- (0,#1+1) node[anchor=east]{$y$};
\foreach \x in {2,4,...,#1}
    {
    \draw[thick] (\x,.2) -- (\x,-.2) node[anchor=north,font=\scriptsize]{$\x$};
    \draw[thick] (-\x,.2) -- (-\x,-.2) node[anchor=north,font=\scriptsize]{$-\x\phm$};
    \draw[thick] (.2,\x) -- (-.2,\x) node[anchor=east,font=\scriptsize]{$\x$};
    \draw[thick] (.2,-\x) -- (-.2,-\x) node[anchor=east,font=\scriptsize]{$-\x$};
    }
}

\newlength{\storea}
\newlength{\storeb}
\newlength{\storec}
\newlength{\storeh}
\newlength{\storek}
\newlength{\storexi}
\newlength{\storexii}
\newlength{\pgwid}
\newlength{\pghit}
\newlength{\pgx}
\newlength{\pgy}

\newcommand{\parabolapoints}[6][2in]{%
\setlength{\storea}{1pt*\real{#2}}
\setlength{\storeb}{1pt*\real{#3}}
\setlength{\storec}{1pt*\real{#4}}
\setlength{\storeh}{1pt*\ratio{\storeb*\real{-1}}{\storea*2}}
\setlength{\storek}{1pt*\ratio{\storeb*\ratio{\storeb}{-1pt}+\storea*\ratio{\storec}{1pt}*4}{\storea*4}}
\setlength{\storexi}{1pt*\real{#5}}
\setlength{\storexii}{1pt*\real{#6}}
\setlength{\pgwid}{\maxof{\maxof{\storexi}{0pt}}{\storexii}-\minof{\minof{\storexi}{0pt}}{\storexii}}
\setlength{\pgwid}{\pgwid*\real{1.2}}
\setlength{\pghit}{\maxof{\maxof{\storec}{0pt}}{\storeh}-\minof{\minof{\storec}{0pt}}{\storeh}}
\setlength{\pghit}{\pghit*\real{1.2}}
\setlength{\pgx}{#1*\ratio{1pt}{\pgwid}}
\setlength{\pgy}{#1*\ratio{1pt}{\pghit}}
\begin{ilpic}[x=\pgx,y=\pgy]

\end{ilpic}
}

%first two arguments are top left and bottom right of graph
%remaining arguments are two points on the graph
\NewDocumentCommand{\linepoints}{>{\SplitArgument{1}{,}}r() >{\SplitArgument{1}{,}}r() >{\SplitArgument{1}{,}}r() >{\SplitArgument{1}{,}}r()} {
	\linepointsSUB #3 #4 #1 #2
}

%first two arguments are top left and bottom right of graph
%third argument is a point on the graph
%fourth argument is the slope
\NewDocumentCommand{\linepointslope}{>{\SplitArgument{1}{,}}r() >{\SplitArgument{1}{,}}r() >{\SplitArgument{1}{,}}r() m}{
	\lineslopeSUB #1 #2 #3{#4}
}

%first two arguments are top left and bottom right of graph
%third argument is the slope
%fourth argument is the y-coordinate of the y-intercept
\NewDocumentCommand{\lineslopeintercept}{>{\SplitArgument{1}{,}}r() >{\SplitArgument{1}{,}}r() m m}{
	\lineslopeSUB #1 #2{0}{#3}{#4}
}

\NewDocumentCommand{\linepointsSUB}{m m m m m m m m} {%
	\lineslopeSUB{#5}{#6}{#7}{#8}{#1}{#2}{(#4-#2)/(#3-#1)}
}

%#1 is x-coordinate of bottom left
%#2 is y-coordinate of bottom left
%#3 is x-coordinate of top right
%#4 is y-coordinate of top right
%#5 is x-coordinate of point
%#6 is y-coordinate of point
%%7 is slope
\NewDocumentCommand{\lineslopeSUB}{m m m m m m m} {%
	\path[name path=grbound]
		(#1-0.9,#2-0.9) rectangle ({#3+0.9},#4+0.9);
	\path[name path=myline]
		(#1-1,{(#7)*(#1-1-#5)+#6}) -- (#3+1,{(#7)*(#3+1-#5)+#6});
	\draw[name intersections = {of = grbound and myline, name = pt},
				thick,<->] (pt-1) -- (pt-2);
}

\newcommand{\GraphFit}[6] {%
	%argument 1 is the dimension of the graph
	%arguments 2 and 3 are the minimum and maximum x-values
	%arguments 4 and 5 are the minimum and maximum y-values
	%argument 6 is the function to graph
	\begin{ilpic}[x={#1/(#3-#2+2)},y={#1/(#5-#4+2)}]
		\draw[step=1,ultra thin] (#2-.5,#4-.5) grid (#3+.5,#5+.5);
		\draw[<->,very thick] (#2-1,0) -- (#3+1,0) node[anchor=south]{$x$};
		\draw[<->,very thick] (0,#4-1) -- (0,#5+1) node[anchor=east]{$y$};
		\foreach \x in {1,...,#3} \draw[thick] (\x,.2) -- (\x,-.2) node[anchor=north,font=\scriptsize]{$\x$};
		\foreach \x in {#2,...,-1} \draw[thick] (\x,.2) -- (\x,-.2) node[anchor=north,font=\scriptsize]{$\x\phm$};
		\foreach \y in {1,...,#5} \draw[thick] (.2,\y) -- (-.2,\y) node[anchor=east,font=\scriptsize] {$\y$};
		\foreach \y in {#4,...,-1} \draw[thick] (.2,\y) -- (-.2,\y) node[anchor=east,font=\scriptsize] {$\y$};
		\clip (#2-.5,#4-.5) rectangle (#3+.5,#5+.5);
		\draw[domain={#2-.5}:{#3+.5}, samples=100,thick] plot (\x,{#6});
	\end{ilpic}
}



\usepackage{calc}
\usepackage{ifthen}
\newcounter{grleft}
\newcounter{grright}
\newcounter{grbot}
\newcounter{grtop}
\newcounter{grwidth}
\newcounter{grheight}
\newcounter{grdiv}
\newcounter{grdif}
\newcommand{\CircleFit}[4][0.4\textwidth] {%
	%argument 1 is the dimension of the graph
	%arguments 2 and 3 are the coordinates of the center
	%argument 4 is the radius
	\setcounter{grleft}{1*\real{#4}-#2}%
	\ifthenelse{\value{grleft}<0}%
	          {\setcounter{grleft}{0}}%
						{\addtocounter{grleft}{1}}%
	\setcounter{grright}{#2+1*\real{#4}}%
	\ifthenelse{\value{grright}<0}%
	          {\setcounter{grright}{0}}%
						{\addtocounter{grright}{1}}%
	\setcounter{grbot}{1*\real{#4}-#3}%
	\ifthenelse{\value{grbot}<1}%
	          {\setcounter{grbot}{0}}%
						{\addtocounter{grbot}{1}}%
	\setcounter{grtop}{#3+1*\real{#4}}%
	\ifthenelse{\value{grtop}<1}%
	          {\setcounter{grtop}{0}}%
						{\addtocounter{grtop}{1}}%
	\setcounter{grwidth}{\value{grleft}+\value{grright}}%
	\setcounter{grheight}{\value{grtop}+\value{grbot}}%
	\ifthenelse{\value{grwidth}<\value{grheight}}{%
		\setcounter{grdif}{\value{grheight}-\value{grwidth}}%
		\ifthenelse{\isodd{\value{grdif}}}{%Set it up so we'll be able to do symmetric stuff
			\addtocounter{grdif}{1}%
			\ifthenelse{\value{grbot}<\value{grtop}}{%
				\addtocounter{grbot}{1}%
			}{%
				\addtocounter{grtop}{1}%
			}%
		}{}%
		\setcounter{grdif}{\value{grdif}*\real{0.5}}%
		\addtocounter{grleft}{\value{grdif}}%
		\addtocounter{grright}{\value{grdif}}%
	}{%
		\setcounter{grdif}{\value{grwidth}-\value{grheight}}%
		\ifthenelse{\isodd{\value{grdif}}}{%Set it up so we'll be able to do symmetric stuff
			\addtocounter{grdif}{1}%
			\ifthenelse{\value{grleft}<\value{grright}}{%
				\addtocounter{grleft}{1}%
			}{%
				\addtocounter{grright}{1}%
			}%
		}{}%
		\setcounter{grdif}{\value{grdif}*\real{0.5}}%
		\addtocounter{grtop}{\value{grdif}}%
		\addtocounter{grbot}{\value{grdif}}%
	}%
	\setcounter{grdiv}{\value{grleft}+\value{grright}+2}%
%	Left: \arabic{grleft}, Right: \arabic{grright}, Bottom: \arabic{grbot}, Top: \arabic{grtop}\\
%	Dif: \arabic{grdif}, Width: \arabic{grwidth}, Height: \arabic{grheight}\\
	\begin{ilpic}[x=#1/\value{grdiv},y=#1/\value{grdiv}]
		\draw[step=1,ultra thin] (-\value{grleft}-.5,-\value{grbot}-.5) grid 
		                         (\value{grright}+.5,\value{grtop}+.5);
		\draw[<->,very thick] (-\value{grleft}-1,0) -- (\value{grright}+1,0) node[anchor=south]{$x$};
		\draw[<->,very thick] (0,-\value{grbot}-1) -- (0,\value{grtop}+1) node[anchor=east]{$y$};
%What the heck--we'll do fussy things where we have dimensions less than one.
		\ifthenelse{\value{grright}>0}{%
			\foreach \x in {1,...,\value{grright}}
				\draw[thick] (\x,.2) -- (\x,-.2) node[anchor=north,font=\scriptsize]{$\x$};
		}{}%If grright=0 do nothing
		\ifthenelse{\value{grleft}>0}{%
			\foreach \x in {1,...,\value{grleft}}
				\draw[thick] (-\x,.2) -- (-\x,-.2) node[anchor=north,font=\scriptsize]{$-\x\phm$};
		}{}%If grleft=0 do nothing
		\ifthenelse{\value{grtop}>0}{%
			\foreach \y in {1,...,\value{grtop}}
				\draw[thick] (.2,\y) -- (-.2,\y) node[anchor=east,font=\scriptsize]{$\y$};
		}{}%If grtop=0 do nothing
		\ifthenelse{\value{grbot}>1}{%
			\foreach \y in {2,...,\value{grbot}}
				\draw[thick] (.2,-\y) -- (-.2,-\y) node[anchor=east,font=\scriptsize]{$-\y$};
		}{}%If grbot=0 do nothing
		\draw[thick] (#2,#3) circle (#4);
		\fill (#2,#3) circle (0.025in);
	\end{ilpic}
}

\newcommand{\LittleGraph}[1]{%
\begin{ilpic}[x=.1in,y=.1in]
\useasboundingbox(-5.9,-6.1) rectangle (5.9,6.1);
\DrawGridSquareBlank 
\clip (-5.9,-5.9) rectangle (5.9,5.9);
#1
\end{ilpic}}

\NewDocumentCommand{\LittleGraphBlank}{O{5} m O{1in}}{%
\begin{ilpic}[x={0.5/(#1+0.9)*#3},y={0.5/(#1+0.9)*#3}]
\draw[<->,thick] (-#1-0.9,0) -- (#1+0.9,0) node[anchor=south]{$x$};
\draw[<->,thick] (0,-#1-0.9) -- (0,#1+0.9) node[anchor=west]{$y$};
#2
\end{ilpic}
}

\NewDocumentCommand{\LittleGraphBlankCompressed}{O{5} m O{1in}}{%
\begin{ilpic}[x={0.5/(#1+0.9)*#3},y={0.5/(#1+0.9)*#3}]
\useasboundingbox(-#1-0.9,-0.5*#3+1ex) rectangle (#1+0.9,0.5*#3-1ex);
\draw[<->,thick] (-#1-0.9,0) -- (#1+0.9,0) node[anchor=south]{$x$};
\draw[<->,thick] (0,-#1-0.9) -- (0,#1+0.9) node[anchor=west]{$y$};
#2
\end{ilpic}
}

\newcommand{\powgraph}[3][1]{%
\draw[<->,very thick,domain={-#2^(1/#3)}:{#2^(1/#3)}] plot (\x,{#1*pow(\x,#3)});
}

\newcommand{\WorkGridSquare}[1][5]{
\begin{ilpic}[x=.2in,y=.2in]\DrawGridSquare[#1]\end{ilpic}
}

\newcommand{\LittleGraphEO}[2][.1in]{%
\begin{ilpic}[x=#1,y=#1]
\useasboundingbox(-5.9,-5.9) rectangle (5.9,5.9);
\DrawGridSquareEO
\clip (-5.9,-5.9) rectangle (5.9,5.9);
#2
\end{ilpic}}

\newcommand{\DrawGridSquareHalf}[1][2]{
\draw[step=1,ultra thin] (-#1-.9,-#1-.9) grid (#1+.9,#1+.9);
\draw[<->,very thick] (-#1-1,0) -- (#1+1,0) node[anchor=south]{$x$};
\draw[<->,very thick] (0,-#1-1) -- (0,#1+1) node[anchor=east]{$y$};
\foreach \x in {1,...,#1}
    {
    \draw[thick] (2*\x,.2) -- (2*\x,-.2) node[anchor=north,font=\scriptsize]{$\x$};
    \draw[thick] (-2*\x,.2) -- (-2*\x,-.2) node[anchor=north,font=\scriptsize]{$-\x\phm$};
    \draw[thick] (.2,2*\x) -- (-.2,2*\x) node[anchor=east,font=\scriptsize]{$\x$};
    \draw[thick] (.2,-2*\x) -- (-.2,-2*\x) node[anchor=east,font=\scriptsize]{$-\x$};
    }
}





\newcommand{\Set}[1]{\bigl\{#1\bigr\}}
\newcommand{\Relat}[1]{\Set{(x,y)\mid#1}}
\newcommand{\SetBuild}[2][x]{\Set{#1\mid #2}}
\newcommand{\Abs}[1]{\left|#1\right|}

\newenvironment{clist}{
\begin{enumerate}
\setlength{\topsep}{0pt}
\setlength{\itemsep}{1pt}
\renewcommand{\theenumi}{\alph{enumi}}
\renewcommand{\labelenumi}{\theenumi.}
}{
\end{enumerate}
}

\newenvironment{Clist}{
\begin{enumerate}
\setlength{\topsep}{0pt}%
\setlength{\itemsep}{1pt}%
\renewcommand{\theenumi}{\Alph{enumi}}
\renewcommand{\labelenumi}{\theenumi.}
}{
\end{enumerate}
}

\newenvironment{Choices}{%
\begin{enumerate}%
\renewcommand{\labelenumii}{\Alph{enumii}.}%
\setlength{\topsep}{0pt}%
\setlength{\itemsep}{1pt}%
}{%
\end{enumerate}%
}

\newcommand{\inv}{^{-1}}
\newcommand{\wordand}{\mbox{ and }}
\newcommand{\wordor}{\mbox{ or }}
\newcommand{\Gp}[1]{\left(#1\right)}

\newcommand{\makerefable}{
\addtocounter{multienum\romannumeral\themultienumdepth}{-1}%
\refstepcounter{multienum\romannumeral\themultienumdepth}}

\newcommand{\mlbl}[1]{\makerefable\label{#1}}

\newcommand{\nref}[1]{\#\ref{#1}}
\newcommand{\nnref}[2]{\#\ref{#1} -- \ref{#2}}

\newcommand{\lrint}[2]{\left(#1,#2\right)}
\newcommand{\Lrint}[2]{\left[#1,#2\right)}
\newcommand{\lRint}[2]{\left(#1,#2\right]}
\newcommand{\LRint}[2]{\left[#1,#2\right]}
\newcommand{\lint}[1]{\left(#1,\infty\right)}
\newcommand{\Lint}[1]{\left[#1,\infty\right)}
\newcommand{\rint}[1]{\left(-\infty,#1\right)}
\newcommand{\Rint}[1]{\left(-\infty,#1\right]}
\newcommand{\eint}{\left(-\infty,\infty\right)}
\newcommand{\neint}[1]{\left(-\infty,#1\right)\cup\left(#1,\infty\right)}

\DeclareMathOperator{\dom}{Dom}
\DeclareMathOperator{\rng}{Rng}

%\newenvironment{amath}{\begin{math}\begin{aligned}}{\end{aligned}\end{math}}
\NewDocumentEnvironment{amath}{b}{$\begin{aligned}#1\end{aligned}$}{}

\NewDocumentCommand{\CS}{m o o}{%
$x^2 #1x+{\IfValueTF{#2}{#2}{\qquad}}%
=%
\IfValueTF{#3}{\Bigl(x+#3\Bigr)^2}{\Bigl(x+\qquad\Bigr)^2}$%
}

\NewDocumentCommand{\CSsquare}{o o}{%
\\\begin{ilpic}[x=1.5em,y=1.5em]
\useasboundingbox(-1,1) rectangle (5,-5);
\draw (0,0) rectangle (5,-5);
\draw (3,0) -- (3,-5);
\draw (0,-3) -- (5,-3);
\draw (1.5,-1.5) node[anchor=mid]{$x^2$};
\IfValueT{#1}{
	\draw (4,-1.5) node[anchor=mid]{$#1 x$};
	\draw (1.5,-4) node[anchor=mid]{$#1 x$};
}
\IfValueT{#2}{
	\draw (4,-4) node[anchor=mid]{$#2$};
}
\end{ilpic}
}

\newenvironment{pw}{\displaystyle\left\{\begin{array}{l@{\quad\mathrm{if}\quad}l}}{\end{array}\right.}


\newlength{\dotsize}
\setlength{\dotsize}{0.025in}

\newcommand{\Cdot}[1]{\filldraw (#1) circle (\dotsize)}
\newcommand{\Odot}[1]{\filldraw[thick,fill=white] (#1) circle (\dotsize)}

\newcommand{\maitemxx}[2]{\mitemxx{$#1$}{$#2$}}
\newcommand{\maitemxxx}[3]{\mitemxxx{$#1$}{$#2$}{$#3$}}
\newcommand{\maitemxxxx}[4]{\mitemxxxx{$#1$}{$#2$}{$#3$}{$#4$}}


\newenvironment{System}{%
	\displaystyle\left\{\begin{array}{l}%
}{%
	\end{array}\right.%
}

\usepackage{array}
\newcolumntype{o}{@{}>{{}}c<{{}}@{}}

\newenvironment{SystemA}{%
	\displaystyle\left\{\begin{array}{rororor}%
}{%
	\end{array}\right.%
}

\newcommand{\N}{\mathbb{N}}
\newcommand{\Z}{\mathbb{Z}}
\newcommand{\Q}{\mathbb{Q}}
\newcommand{\R}{\mathbb{R}}

\newcommand{\remainder}{\;\textsc{r}}

\newcommand{\nlgraph}[2][6]{%
\begin{tikzpicture}[x=0.2in,y=0.2in,baseline=-0.2*0.2in]
\draw[<->,semithick] (-#1-1,0) -- (#1+1,0);
\foreach \x in {0,...,#1} {
	\draw (\x,0.2) -- (\x,-0.2);
	\draw (\x,-0.7) node[anchor=base]{\footnotesize{$\x$}};
}
\foreach \x in {1,...,#1} {
	\draw (-\x,0.2) -- (-\x,-0.2);
	\draw (-\x,-0.7) node[anchor=base,xshift=-0.75ex]{\footnotesize{$-\x$}};
}
#2
\end{tikzpicture}
}

\usepackage{ifthen}


\NewDocumentCommand{\nlgraphgen}{O{-} m m m m O{0}}{%
	\draw[ultra thick,#1] (#2,#6) -- (#3,#6);
	\ifthenelse{\equal{#1}{<-}}{}{\filldraw[thick,black,fill=#4] (#2,#6) circle (0.125*0.2in);}
	\ifthenelse{\equal{#1}{->}}{}{\filldraw[thick,black,fill=#5] (#3,#6) circle (0.125*0.2in);}
}
\NewDocumentCommand{\nlgraphoo}{m m O{0}}{%
	\nlgraphgen{#1}{#2}{white}{white}[#3]
}
\NewDocumentCommand{\nlgraphoc}{m m O{0}}{%
	\nlgraphgen{#1}{#2}{white}{black}[#3]
}
\NewDocumentCommand{\nlgraphco}{m m O{0}}{%
	\nlgraphgen{#1}{#2}{black}{white}[#3]
}
\NewDocumentCommand{\nlgraphcc}{m m O{0}}{%
	\nlgraphgen{#1}{#2}{black}{black}[#3]
}
\NewDocumentCommand{\nlgraphao}{O{6} m O{0}}{%
	\nlgraphgen[<-]{-#1-1}{#2}{white}{white}[#3]
}
\NewDocumentCommand{\nlgraphac}{O{6} m O{0}}{%
	\nlgraphgen[<-]{-#1-1}{#2}{black}{black}[#3]
}
\NewDocumentCommand{\nlgraphoa}{O{6} m O{0}}{%
	\nlgraphgen[->]{#2}{#1+1}{white}{white}[#3]
}
\NewDocumentCommand{\nlgraphca}{O{6} m O{0}}{%
	\nlgraphgen[->]{#2}{#1+1}{black}{black}[#3]
}

\newcounter{argcount}
\NewDocumentCommand{\nlrough}{O{2in} >{\SplitArgument{3}{,}} m m}{%
\countargs#2
\nlroughpro[#1]#2{#3}
}

\NewDocumentCommand{\nlroughpro}{O{2in} m m m m m}{%
%\setcounter{argcount}{0}%
%\IfValueT{#2}{\addtocounter{argcount}{1}}%
%\IfValueT{#3}{\addtocounter{argcount}{1}}%
%\IfValueT{#4}{\addtocounter{argcount}{1}}%
%\IfValueT{#5}{\addtocounter{argcount}{1}}%
\begin{tikzpicture}[x={#1/(\value{argcount}+1)},y=0.2in,baseline=-0.2*0.2in]
\draw[semithick,<->] (0,0) -- (\value{argcount}+1,0);
\foreach \x in {1,...,\value{argcount}} \draw (\x,0.2) -- (\x,-0.2);
\IfValueT{#2}{\draw (1,-0.8) node[anchor=base]{\footnotesize$#2$};}
\IfValueT{#3}{\draw (2,-0.8) node[anchor=base]{\footnotesize$#3$};}
\IfValueT{#4}{\draw (3,-0.8) node[anchor=base]{\footnotesize$#4$};}
\IfValueT{#5}{\draw (4,-0.8) node[anchor=base]{\footnotesize$#5$};}
#6
\end{tikzpicture}%
}

\newcommand{\atype}{}
\newcommand{\lend}{}
\newcommand{\rend}{}
\newcommand{\lendv}{}
\newcommand{\rendv}{}
\newcommand{\nla}{}
\newcommand{\nlb}{}

%#1 is the size of the number line
%#2 is the values to show (up to 3 of them)
%#3 is the types of endpoints for the first graph:
%  a,o, or c; exactly 2
%#4 is the endpoints for the first graph:
%  one if there is an arrow, two otherwise
%#5 is as #3
%#6 is as #4
\NewDocumentCommand{\twonlgraph}{O{2in} >{\SplitArgument{3}{,}}m >{\SplitArgument{1}{,}}m >{\SplitArgument{1}{,}}m >{\SplitArgument{1}{,}}m >{\SplitArgument{1}{,}}m}%
{%
	\countargs#2%
	\nlroughpro[#1]#2{
		\parsetype#3#4%
		\nlgraphgen[\atype]{\lendv}{\rendv}{\lend}{\rend}[0.5]%
		\parsetype#5#6%
		\nlgraphgen[\atype]{\lendv}{\rendv}{\lend}{\rend}[0.25]%
	}
}
\NewDocumentCommand{\countargs}{m m m m}{%
	\setcounter{argcount}{0}%
	\IfValueT{#1}{\addtocounter{argcount}{1}}%
	\IfValueT{#2}{\addtocounter{argcount}{1}}%
	\IfValueT{#3}{\addtocounter{argcount}{1}}%
	\IfValueT{#4}{\addtocounter{argcount}{1}}%
}
\NewDocumentCommand{\parsetype}{m m m m}{%
	\ifthenelse{\equal{#1}{a}}{%
		\renewcommand{\atype}{<-}%
		\renewcommand{\lendv}{0}%
		\renewcommand{\rendv}{#3}%
	}{%
		\renewcommand{\atype}{-}%
		\renewcommand{\lendv}{#3}%
	}
	\ifthenelse{\equal{#1}{o}}{%
		\renewcommand{\lend}{white}%
	}{%
		\renewcommand{\lend}{black}%
	}%
	\ifthenelse{\equal{#2}{a}}{%
		\renewcommand{\atype}{->}%
		\renewcommand{\rendv}{\value{argcount}+1}%
	}{%
		\IfValueT{#4}{\renewcommand{\rendv}{#4}}%
	}%
	\ifthenelse{\equal{#2}{o}}{%
		\renewcommand{\rend}{white}%
	}{%
		\renewcommand{\rend}{black}%
	}%
}

\newcommand{\twopartans}[2]{%
\begin{nmenum}%
\mitemxx{#1}{#2}%
\end{nmenum}%
}
\newcommand{\threepartans}[3]{%
\begin{nmenum}%
\mitemxxx{#1}{#2}{#3}%
\end{nmenum}%
}
\newcommand{\threepartansb}[3]{%
\begin{nmenum}%
\mitemxxox{#1}{#2}{#3}%
\end{nmenum}%
}

\newcommand{\withsketch}[2]{%
Sketch: #1\qquad #2
}

\newcommand{\fourpairs}[2]{\tab$(#1,#2)$\tab$(-#1,#2)$\tab$(#1,-#2)$\tab$(-#1,-#2)$}
\newcommand{\fourpairsb}[4]{\tab$(#1)$\tab$(#2)$\tab$(#3)$\tab$(#4)$}

\newcommand{\mg}[1]{#1\;\mathrm{mg}}

\newcommand{\graphpage}{\clearpage%
\noindent\begin{tikzpicture}[x=0.25in,y=0.25in]
\useasboundingbox (-0.5\textwidth,-0.5\textheight) rectangle (0.5\textwidth,0.5\textheight);
\draw[step=1,black!50!white] (-0.5\textwidth,-0.5\textheight) grid (0.5\textwidth,0.5\textheight);
\end{tikzpicture}}


\NewDocumentCommand{\threesln}{>{\SplitArgument{2}{,}} O{x,y,z} >{\SplitArgument{2}{,}} m}{%
	\threeslnSUB #1 #2%
}

\NewDocumentCommand{\threeslnSUB}{m m m m m m}{%
	\ensuremath{#1=#4,\ #2=#5,\ #3=#6}%
}